\BeleskeID{09_1-s052}
% element: 09_1-s052:n01
U ovoj oblasti možemo jednostavno izraziti izlaz preko ulaza. Polazna relacija, koja se ponavlja u izvođenju, jeste
\[k_{n,L}(V_{DD}-V_O-V_{Tn,L})^2=k_{n,D}(V_I-V_{Tn,D})^2\]
Za fizičku granu uzimamo nenegativne napone prezasićenja. Kvadratni koren zatim vodi do linearne zavisnosti izlaza od ulaza. Zbog toga izvod ne zavisi od radne tačke u toj oblasti. Za razmatrano logičko kolo potrebno je da apsolutna vrednost pojačanja bude veća od jedan, što određuje odnos parametara pogonskog tranzistora i opterećenja.
